\documentclass{beamer}					% Document class

\usepackage[english]{babel}				% Set language
\usepackage[utf8]{inputenc}			% Set encoding

\mode<presentation>						% Set options
{
  \usetheme{default}					% Set theme
  \usecolortheme{default} 				% Set colors
  \usefonttheme{default}  				% Set font theme
  \setbeamertemplate{caption}[numbered]	% Set caption to be numbered
  \setbeamercolor{footnote}{fg=gray}
  \setbeamerfont{footnote}{size=\footnotesize}
}

% Uncomment this to have the outline at the beginning of each section highlighted.
\AtBeginSection[]
{
 \begin{frame}{Outline}
   \tableofcontents[currentsection]
 \end{frame}
}


\usepackage{graphicx}					% For including figures
\usepackage{booktabs}					% For table rules
\usepackage{hyperref}					% For cross-referencing

\usepackage{tikz}                     % For drawing diagrams
\usepackage{caption} 
\usepackage[normalem]{ulem}

\captionsetup[figure]{font={small, sf},labelformat=empty} % supprimer "Figure 1:" dans les captions des figures
\captionsetup[table]{font={small, sf},labelformat=empty} 

\usetikzlibrary{calc}

\title{Speech to emotion classification}	% Presentation title
\subtitle{Large-scale Generative Models for NLP and Speech Processing - Project}		% Presentation subtitle
\author{Maloe Aymonier, Baptiste Geisenberger, Astrid Giuliani}								% Presentation author
%\institute{Télécom Paris}					% Author affiliation
\date{25/03/2026}									% Today's date	

\begin{document}

% Title page
% This page includes the informations defined earlier including title, author/s, affiliation/s and the date
\begin{frame}
  \titlepage
\end{frame}

% Outline
% This page includes the outline (Table of content) of the presentation. All sections and subsections will appear in the outline by default.
% \begin{frame}{Outline}
%   \tableofcontents
% \end{frame}

% The following is the most frequently used slide types in beamer
% The slide structure is as follows:
%
%\begin{frame}{<slide-title>}
%	<content>
%\end{frame}



\section{Task definition} 
% Pick a task and dataset

\section{Dataset}
% Analyze the data

\section{Methodology}
% Build a model

\section{Results}
% Run experiments, show results


\section{Discussion}
% Analyze the errors made by the model

\end{document}
